\section{Convexity}
\index{concave}%
\index{convex}%

\begin{definition}[convex function]
\mymark{**}
Let $I\subset \RR$ be an interval.
A function $f\colon I\to \RR$
is said to be
\emph{convex}
\mymargin{convex function}%
\index{function!convex}%
\index{convex!function}%
if for every $x,y\in I$ and for every $t\in [0,1]$ we have
\[
f((1-t)x + ty) \le (1-t) f(x) + t f(y).
\]

Similarly, we say that $f$ is \emph{concave} 
\mymargin{concave function}%
\index{function!concave}%
\index{concave!function}%
if the reverse inequality holds:
\[
f((1-t)x + ty) \ge (1-t) f(x) + t f(y)
\]
(or, equivalently, if $-f$ is convex).
\end{definition}

Observe that the line in the plane passing through the points $(x,f(x))$ and
$(y,f(y))$ can be parametrized uniformly for $t\in \RR$
by
\[
  (1-t) (x,f(x)) + t(y,f(y)) = ((1-t)x + ty, (1-t) f(x) + tf(y)).
\]
Clearly, for $t=0$ we obtain the point $(x,f(x))$, for $t=1$ the point
$(y,f(y))$, and for $t\in[0,1]$ the segment joining these points. The convexity
condition of the function $f$ thus corresponds to requiring that every chord
(segment) joining two points of the graph lies "above" the graph of the
function.

\begin{definition}[convex set]
\mymark{*}
A set $E\subset \RR^n$ is said to be \emph{convex}%
\mymargin{convex}%
\index{convex} if given
any two points $a,b\in E$, the entire segment $[a,b]=\ENCLOSE{(1-t)a+tb\colon
t\in [0,1]}$ is contained in $E$.
\end{definition}

\begin{theorem}[epigraph of convex functions]
Let $I\subset \RR$ and $f\colon I\subset \RR\to \RR$ be a function.
Then the following are equivalent:
\begin{enumerate}
\item $I$ is an interval and $f$ is convex;
\item the \emph{epigraph of $f$}%
\mymargin{epigraph of $f$}%
\index{epigraph}%
(or \emph{supergraph})
\index{epigraph}%
i.e., the set
\[
  E = \ENCLOSE{(x,y)\in \RR^2\colon x\in I, y\ge f(x)}
\]
is convex.
\end{enumerate}

For concave functions, the \emph{subgraph} $\ENCLOSE{(x,y)\colon y\le f(x)}$
will be convex.
\end{theorem}
%
\begin{proof}
Suppose that $I$ is an interval and $f$ is convex. 
To prove that the epigraph $E$ is convex, consider two points $a,b\in E$ 
and any point $p$ on the segment $[a,b]$.
If $a=(x_a, y_a)$, $b=(x_b,y_b)$, $p=(x_p, y_p)$
then there exists a $t\in [0,1]$ such that $x_p = (1-t)x_a + t x_b$ and
$y_p=(1-t)y_a + t y_b$.
Since $a,b\in E$, we know that $y_a \ge f(x_a)$ and $y_b\ge f(x_b)$. 
Thus, necessarily, we have
\[
  y_p \ge (1-t)f(x_a) + t f(x_b).
\]
But since $f$ is convex, we have:
\[
  (1-t)f(x_a) + t f(x_b) \ge f((1-t)x_a + t x_b) = f(x_p).
\]
Thus $y_p\ge f(x_p)$, which means $p\in E$.

Conversely, suppose we know that $E$ is convex. Let $x,y\in I$ be any points.
Then the points $a=(x,f(x))$ and $b=(y,f(y))$ are certainly points of $E$, and
therefore the entire segment $[a,b]$ must be contained in $E$. Thus, for every
$t\in [0,1]$, the point $p = ((1-t)x + t y,$ $(1-t)f(x)+ tf(y))$ must lie in
$E$. Firstly, it must be that $(1-t)x+ty\in I$, and if this is true for every
$t\in[0,1]$, it means that $I$ is an interval. Secondly, if $p\in E$, it means
that
\[
  (1-t)f(x) +t f(y) \ge f((1-t)x + t y)
\]
which corresponds to the definition of a convex function.
\end{proof}


\begin{lemma}[incremental ratio of a convex function]
\mymark{*}%
\label{lemma:547091}%
Let $I$ be an interval of $\RR$ and let $f\colon I\to \RR$.
Given $x,y\in I$ with $x\neq y$, we define the \emph{incremental ratio}
of $f$ as:
\[
  R(x,y) = \frac{f(y) - f(x)}{y-x}.
\]
Then the following conditions are equivalent:
\begin{enumerate}
\item $f$ is convex;
\item for every $x,y,z\in I$ if $x<y<z$ then $R(x,y)\le R(y,z)$;
\item for every $x,y,z\in I$ if $x<y<z$ then $R(x,y)\le R(x,z)$;
\item for every $x,y,z\in I$ if $x<y<z$ then $R(x,z)\le R(y,z)$;
\item the function $R(x,y)$ is increasing in each of the two variables.
\end{enumerate}
\end{lemma}
%
\begin{proof}
Attention:
the lemma is obvious if one uses the correct geometric interpretation
(the incremental ratio is the slope of the corresponding chord).
What follows is the algebraic translation of what
is geometrically obvious but inevitably becomes cumbersome
and more difficult to understand.

Let $x,y,z\in I$ with $x<y<z$.
Setting $t=(y-x)/(z-x)$, we have $y=(1-t)x + tz$,
 $y-x = t(z-x)$, $z-y = (1-t)(z-x)$.
Then we have
 \begin{equation*}
 \begin{aligned}
 R(x,z) - R(x,y)
 &= \frac{f(z)-f(x)}{z-x} - \frac{f(y)-f(x)}{y-x} \\
  &= t\frac{f(z)-f(x)}{y-x} - \frac{f(y)-f(x)}{y-x} \\
  &= \frac{tf(z) + (1-t) f(x) - f(y)}{y-x}
 \end{aligned}
 \end{equation*}

The convexity condition of $f$ is
\[
  f(y) \le (1-t)f(x) + tf(z)
\]
and is therefore equivalent to the condition $R(x,y) \le R(x,z)$.
Thus, conditions 1 and 3 are equivalent.

But with a direct verification, we observe that
\[
  R(x,z) = t R(x,y) + (1-t) R(y,z)
\]
from which we obtain
\[
  R(y,z) - R(x,z) = t[R(y,z) - R(x,y)]
\]
or also
\[
 R(x,z) - R(x,y) = (1-t) [R(y,z) - R(x,y)].
\]
It follows that the quantities
\[
  R(y,z) - R(x,y), \qquad
  R(x,z) - R(x,y), \qquad
  R(y,z) - R(x,z)
\]
all have the same sign. Therefore, conditions 2, 3, and 4 are equivalent (if one
of the three holds, all three hold).

If the three conditions 2, 3, and 4 hold, it is easy to verify that the function
$R(x,y)$ is increasing in both variables. First of all, by symmetry, since
$R(x,y) = R(y,x)$, it is sufficient to verify that $R(x,y)$ is increasing in the
second variable $y$ for each fixed $x$. Therefore, given $z>y$, we need to show
that $R(x,z) \ge R(x,y)$.
We then have three possibilities depending on whether $x<y$ or $y<x<z$ or $z<x$.
In the first case, we have $x<y<z$, and thus the inequality $R(x,y) \le R(x,z)$
corresponds to condition 3.
In the second case, we have $y<x<z$, and the condition $R(x,y)\le R(x,z)$ can be
written as $R(y,x) \le R(x,z)$, which, by appropriately rearranging the
variables, is condition 2. If, finally, $y < z < x$, the condition $R(x,y) \le
R(x,z)$ can be written as $R(y,x) \le R(z,x)$, which, by rearranging the
variables, is condition 4.

Conversely (and finally), if the function $R(x,y)$ is increasing in both
variables, in particular, it is increasing in the second variable, and therefore
if $x<y<z$, we have $R(x,y) \le R(x,z)$. It follows that condition 5 implies 3
and therefore all other conditions.
\end{proof}

\begin{theorem}
\mymark{***}
Let $I\subset \RR$ be an interval and $f\colon I \to \RR$ a differentiable
function on all of $I$.
Then the following are equivalent:
\begin{enumerate}
\item $f$ is convex;
\item for every $x_0 \in I$ and for every $x\in I$, we have
\[
   f(x) \ge f'(x_0) (x-x_0) + f(x_0)
\]
(geometrically: the graph of the function lies above the tangent line);
\item $f'$ is increasing.
\end{enumerate}

Similarly, for concave functions, the graph will "lie below" the tangent line,
and the derivative will be decreasing.
\end{theorem}
%
\begin{proof}
\mymark{**}
Observe that
\[
  f'(x_0) = \lim_{x\to x_0} R(x_0,x).
\]
If $f$ is convex, then, by the lemma, the incremental ratio $R(x_0,x)$ is
increasing, and therefore $f'(x_0) = \inf_{x>x_0} R(x_0,x)$. In particular,
$f'(x_0) \le R(x_0,x)$ for every $x> x_0$. Similarly, we find $f'(x_0) \ge
R(x_0,x)$ if $x<x_0$.
In any case, it follows that for every $x$, we have
\[
(R(x_0,x)- f'(x_0))(x-x_0)\ge 0
\]
or equivalently
\[
  f(x) - f(x_0) - f'(x_0)(x-x_0) \ge 0.
\]
Thus, condition 1 implies condition 2.

If condition 2 holds, given $x,y \in I$, we have
\[
  f(x) - f(y) \ge f'(y)(x-y)
\]
if we exchange $x$ and $y$ and change the sign of both sides, we obtain
\[
  f(x) - f(y) \le f'(x)(x-y)
\]
putting together the two inequalities,
if we now assume that $x>y$, we obtain precisely
$f'(x) \ge f'(y)$, i.e., $f'$ is increasing (condition 3).

Now suppose we know that $f'$ is increasing 
and prove that if $x<y<z$, then $R(x,y) \le R(y,z)$.
By the mean value theorem, there exist $s\in (x,y)$ 
and $t\in (y,z)$ such that $f'(s) = R(x,y)$ and $f'(t) = R(y,z)$.
But $s<t$, thus $f'(s) \le f'(t)$, i.e., $R(x,y) \le R(y,z)$
as we wanted to prove.
\end{proof}

\begin{corollary}[convexity criterion via second derivative]
\mymark{***}
Let $I\subset \RR$ be an interval and let $f\colon I \to \RR$ be a twice
differentiable function (i.e., $f$ is differentiable and $f'$ is also
differentiable).
Then $f$ is convex if and only if $f''(x)\ge 0$ for every $x\in I$.
Similarly, $f$ is concave if and only if $f''\le 0$.
\end{corollary}
\begin{proof}
\mymark{***}
By the previous criterion, $f$ is convex if and only if $f'$ is increasing. By
the monotonicity criterion, $f'$ is increasing if and only if $f'' \ge 0$.
Similar considerations apply to concavity.
\end{proof}

\begin{theorem}
Let $a\in \RR$, $b\in \bar \RR$, $a<b$.
Let $f\colon [a,b)\to \RR$ be a convex function 
in $(a,b)$ and continuous at $a$. 
Then $f$ is convex on the entire $[a,b)$. 
A similar result holds for functions defined on 
left-open intervals $(a,b]$
or on intervals open on both sides $(a,b)$.
\end{theorem}
%
\begin{proof}
It suffices to show that if $x<y<z$, then $R(x,y)\le R(y,z)$.
The inequality holds for $x>a$ since $f$ is convex 
in $(a,b)$, and thus, by taking the limit as $x\to a$, 
thanks to the continuity of $f$ at $a$,
the inequality is also proven when $x=a$.
\end{proof}

\begin{example}
The function $f(x) = \sqrt{x}$ is defined on $[0,+\infty)$ but is differentiable
only on $(0,+\infty)$. Its derivative is $f'(x) = x^{-\frac 1 2 }/2$ and the
second derivative is $f''(x) = -x^{-\frac 3 2}/4 < 0$. Thus, the function is
concave on the open interval $(0,+\infty)$. But being continuous, we can
conclude that $f$ is concave on the entire domain $[0,+\infty)$.
\end{example}

\begin{comment}
\begin{theorem}[continuity of convex functions]
Let $a,b \in \bar \RR$ with $a< b$.
Let $f\colon (a,b) \to \RR$ be a convex function. Then $f$ is continuous.
\end{theorem}
%
\begin{proof}
Let $x_0 \in (a,b)$ and let $y,z \in (a,b)$ with $y < x_0 < z$.
By the lemma on incremental ratios, we know that for every $x\in (y,z)$ we have
\[
   R(x_0, y) \le R(x_0,x) \le R(x_0,z).
\]
In particular, there exists a constant $C$ such that
\[
  \abs{R(x_0,x)} \le C,\qquad \forall x \in (y,z).
\]
Multiplying by $\abs{x-x_0}$, we then obtain
\[
   \abs{f(x) - f(x_0)} \le C \abs{x-x_0}
\]
and as $x\to x_0$, the right side tends to zero, and therefore by comparison,
the left side must also tend to zero. Thus $f(x)\to f(x_0)$
and $f$ is continuous at $x_0$.
\end{proof}
\end{comment}

\begin{theorem}[differentiability of convex functions]
    Let $I\subset \RR$ be an interval and let $f\colon I \to \RR$ be a convex
  function.
    If $x_0\in (\inf I, \sup I)$, the right and left derivatives of $f$ at $x_0$
  exist and are finite:
  \begin{align*}
    f'_+(x_0) &= \lim_{x\to x_0^+} \frac{f(x)-f(x_0)}{x-x_0}\\
    f'_-(x_0)  &= \lim_{x\to x_0^-} \frac{f(x)-f(x_0)}{x-x_0}  
  \end{align*}
  and we have
  \[
    f'_-(x_0) \le f'_+(x_0).
  \]
  Moreover, if $\inf I < x_1 < x_2 < \sup I$, we have 
  \[
    f'_+(x_1) \le f'_-(x_2).
  \]
  It follows that $f$ is continuous at all interior points of $I$ 
  and that the set of points where $f$ is not differentiable 
  is at most countable.
  \end{theorem}
  %
  \begin{proof}
    In lemma~\ref{lemma:547091}, we observed that the incremental ratio
    \[
      R(x_0,x) = \frac{f(x)-f(x_0)}{x-x_0}
    \] 
    of a convex function is increasing with respect to the variable $x$. 
    Thus, by theorem~\ref{th:limite_funzione_monotona}, we 
    have 
    \[
      f'_-(x_0) = \sup_{x<x_0} R(x_0,x)>-\infty, \qquad 
      f'_+(x_0) = \inf_{x>x_0} R(x_0,x)<+\infty
    \]
    and since if $x<x_0$ and $y>x_0$, we have $R(x_0,x)\le R(x_0,y)$, 
    we deduce that $f'_-(x_0)\le f'_+(x_0)$, and both limits 
    are thus finite. If $x_1 < x_2$, then for any 
    point $x$ with $x_1<x<x_2$, we have (again by lemma~\ref{lemma:547091})
    \[
      R(x_1,x) \le R(x,x_2)
    \]
    and thus, taking the limit, we find $f'_+(x_1) \le f'_-(x_2)$.
    The right and left derivatives thus exist at all interior points 
    of the interval $I$. 
        In particular, $f$ is continuous at such points since it is continuous
    both from the right
    and from the left.
    The points of non-differentiability of $f$ are the points where 
    the right and left derivatives differ. 
    But to each point $x$ of non-differentiability, 
        I can associate the non-empty open interval $I_x = (f'_-(x),f'_+(x))$,
    and for
        the properties just seen, it is clear that these intervals are pairwise
    disjoint
        since if $x_1<x_2$, the right endpoint of $I_{x_1}$ is less than or
    equal to the left endpoint
    of $I_{x_2}$.
        Since each of these intervals contains at least one rational number, we
    conclude that
        the number of points of discontinuity cannot be greater than the
    cardinality
    of the rational numbers.
  \end{proof}
  
  \begin{theorem}[supporting line]
    \label{th:supporto_convessa}%
    Let $I\subset \RR$ be an interval, $f\colon I\to\RR$ a convex function, 
        and $x_0$ an interior point of $I$. Then there exists an affine linear
    function
    \[
      L(x) = mx+q
    \]
    such that $f(x_0) = L(x_0)$ and for every $x\in I$, we have $f(x)\ge L(x)$.
  \end{theorem}
  %
  \begin{proof}
        It suffices to choose any $m\in [f'_-(x_0),f'_+(x_0)]$ and consider the
    affine linear function $L(x) = m(x-x_0) + f(x_0)$. If $x>x_0$, then we have
    $R(x_0,x)\ge m$,
    and thus $f(x) - f(x_0) \ge m (x-x_0) = L(x)$. Conversely, if $x<x_0$, 
    we have $R(x_0,x)\le m$, hence, again, $f(x)-f(x_0) \ge m(x-x_0) = L(x)$.
  \end{proof}
  %
  The following theorem can also be stated for integrals 
  as we will see in theorem~\ref{th:jensen}. 
  The proof is essentially identical and is valid even 
  for convex functions of multiple variables.
  %
  \begin{theorem}[barycentric combinations/Jensen's inequality]
    \mymark{*}%
    \label{th:combinazioni_baricentriche}%
    \index{barycentric combinations}%
    \index{inequality!Jensen's}%
    \index{Jensen!inequality}%
        If $f$ is a convex function defined on an interval $I$, given $x_1,
    \dots, x_n \in I$ and $\lambda_1, \dots, \lambda_n\in \RR$ such that
    $\sum_{k=1}^n \lambda_k = 1$ and $\lambda_k \ge 0$ for each $k=1, \dots, n$,
    then
    \[
      f\enclose{\sum_{k=1}^n \lambda_k x_k}
      \le \sum_{k=1}^n \lambda_k f(x_k).
    \]
    For concave functions, the reverse inequality holds.
    \end{theorem}
    %
    \begin{comment}
    \begin{proof}
        We proceed by induction on $n$. In the case $n=1$, we have
    $\lambda_1=1$, and the two sides of the inequality are indeed equal.
    Assuming the theorem is proven for a certain $n$, we proceed to prove it for
    $n+1$.
    Observe that
    \begin{align*}
      \sum_{k=1}^{n+1} \lambda_k x_k
      &= \sum_{k=1}^n \lambda_k x_k  + \lambda_{n+1} x_{n+1} \\
            &= (1-\lambda_{n+1})\sum_{k=1}^n\frac{\lambda_k}{1-\lambda_{n+1}}
      x_k + \lambda_{n+1} x_{n+1}.
    \end{align*}
    Since
    \[
      \sum_{k=1}^{n+1} \lambda_k = 1
    \]
    we have
    \[
      \sum_{k=1}^n \frac{\lambda_k}{1-\lambda_{n+1}}
      = \frac{1-\lambda_{n+1}}{1-\lambda_{n+1}} = 1.
    \]
    Thus, by the inductive hypothesis, we have
    \[
      f\enclose{\sum_{k=1}^n\frac{\lambda_k}{1-\lambda_{n+1}} x_k}
      \le \sum_{k=1}^n \frac{\lambda_k}{1-\lambda_{n+1}}f(x_k).
    \]
    Using the convexity of $f$ again with $t=\lambda_{n+1}$, we have
    \begin{align*}
    f\enclose{\sum_{k=1}^{n+1} \lambda_k x_k}
        &=f\enclose{(1-\lambda_{n+1})\sum_{k=1}^n
    \frac{\lambda_k}{1-\lambda_{n+1}}x_k + \lambda_{n+1}x_{n+1}}\\
        &\le (1-\lambda_{n+1})f\enclose{\sum_{k=1}^n
    \frac{\lambda_k}{1-\lambda_{n+1}}x_k} + \lambda_{n+1}f(x_{n+1}) \\
        &\le (1-\lambda_{n+1})\sum_{k=1}^n \frac{\lambda_k}{1-\lambda_{n+1}}
    f(x_k) + \lambda_{n+1} f(x_{n+1})\\
    &= \sum_{k=1}^{n+1}\lambda_k f(x_k).
    \end{align*}
    as we wanted to prove.
    \end{proof}
  \end{comment}
  %    
  \begin{proof}
  Let 
  \[
    \bar x =  \sum_{k=1}^n \lambda_k x_k.
  \]
  By the previous theorem~\ref{th:supporto_convessa}, 
  we know that there exists a line $L(x) = mx + q$ 
  such that $L(\bar x) = f(\bar x)$ 
  and $L(x)\le f(x)$ for every $x \in I$.
  Then we have 
  \[
  L \enclose{\sum_{k=1}^n \lambda_k x_k}
  = \sum_{k=1}^n m \lambda_k x_k + q \sum_{k=1}^n \lambda_k 
  = \sum_{k=1}^n \lambda_k L(x_k).  
  \]
  Thus 
  \[
  f(\bar x) = L(\bar x) = \sum_{k=1}^n \lambda_k L(x_k) \le 
  \sum_{k=1}^n \lambda_k f(x_k).
  \]
  \end{proof}

\begin{example}[inequality between arithmetic mean and geometric mean]
\label{ex:AMGM}%
  \mymark{*}%
\index{inequality!arithmetic mean geometric mean}%
Observe that the function $f(x) = \ln x$ is concave, indeed we have
$f''(x) = -1/x^2 < 0$. Thus, by the previous theorem, if $\lambda_1 + \dots +
\lambda_n =1$, $\lambda_k \ge 0$, we have
\[
    \ln\enclose{\sum_{k=1}^n \lambda_k x_k}
    \ge \sum_{k=1}^n \lambda_k \ln x_k.
\]
Exponentiating both sides, we obtain
\[
  \sum_{k=1}^n \lambda_k x_k \ge \prod_{k=1}^n x_k^{\lambda_k}.
\]
In the particular case $\lambda_k = 1/n$, we obtain
the inequality between \emph{arithmetic mean} (AM for anglophones)
and \emph{geometric mean}
\mymargin{arithmetic geometric mean}%
\index{arithmetic geometric mean}%
\index{mean!arithmetic}%
\index{mean!geometric}%
\index{AM}%
\index{GM}%
(GM):
\[
  \frac{x_1 + \dots + x_n}{n} \ge \sqrt[n]{x_1\cdots x_n}.
\]
\end{example}

\begin{exercise}[subadditivity of concave functions]
\index{subadditivity!of concave functions}%
\index{function!concave!subadditive}%
\index{concave!subadditive}%
Let $f\colon [0,+\infty) \to \RR$ be a concave function with $f(0)\ge 0$. Then
$f$ is subadditive, i.e.:
\[
  f(x+y) \le f(x) + f(y),\qquad \forall x,y\ge 0.
\]
\end{exercise}
%
\begin{proof}
If $x=y=0$, the inequality is obvious.
Otherwise, $x+y>0$, and we have
\begin{align*}
f(x) &= f\enclose{\frac{y}{x+y}\cdot 0 + \frac{x}{x+y}\cdot (x+y)}\\ &\ge
\frac{y}{x+y}f(0) + \frac{x}{x+y}f(x+y)
\ge \frac{x}{x+y}f(x+y).
\end{align*}
Exchanging $x$ with $y$ and summing, we obtain:
\[
  f(x) + f(y) \ge \frac{x}{x+y}f(x+y) + \frac{y}{x+y}f(x+y) = f(x+y).
\]
\end{proof}

Note that by expanding the inequality $(x-y)^2\ge 0$, we obtain:
\[
  x y \le \frac{x^2 + y^2}{2}.
\]
This inequality can be generalized to different exponents 
as seen in the following.

\begin{theorem}[Young's inequality]%
\label{th:young}%
\index{Young!inequality}%
\index{inequality!Young's}%
If $p>1$, $q>1$, $\frac{1}{p}+\frac{1}{q}=1$, $x,y\ge 0$:
\begin{equation}\label{eq:young}
  x y \le \frac{x^p}{p} + \frac{y^q}{q}.
\end{equation}
\end{theorem}
%
\begin{proof}
We use the convexity inequality for the function $\exp(x)=e^x$. 
Since the second derivative of the exponential is positive, the function $\exp$ 
is convex, and thus, if $x>0$ and $y>0$, we have:
\begin{align*}
 x\cdot y 
  &= \exp(\ln x + \ln y)
  = \exp\enclose{\frac 1 p \ln (x^p) + \frac 1 q \ln (y^q)}  \\
  &\le \frac 1 p \cdot \exp \ln (x^p) + \frac 1 q \exp \ln (y^q)
  = \frac{x^p}{p} + \frac{y^q}{q}.
\end{align*}
If $x=0$ or $y=0$, the inequality is obvious since the left side 
of~\eqref{eq:young} is zero.
\end{proof}

\begin{theorem}[Hölder's inequality]
If $a_k\ge 0$ and $b_k\ge 0$ for $k\in \NN$
and let $p>1$ and $q>1$ such that $\frac 1 p + \frac 1 q = 1$.
Then 
\[
 \sum_k a_k b_k  \le \enclose{\sum_k a_k^p}^{1 \over p} \cdot \enclose{\sum_k b_k^q}^{1\over q}.
\]
\end{theorem}
\begin{proof}
Let 
\[
A = \enclose{\sum_k a_k^p}^{1 \over p}, \qquad 
B = \enclose{\sum_k b_k^q}^{1\over q}.
\]
If $A>0$ and $B>0$, we can divide the left side of the inequality 
by the right side and apply Young's inequality~\eqref{eq:young}:
\begin{align*}
  \frac{\sum_k a_k b_k}{A\cdot B} 
  & = \sum_k \frac{a_k}{A}\cdot \frac{b_k}{B} 
  \le \sum_k \frac 1 p \frac{a_k^p}{A^p} + \frac 1 q \frac{b_k^q}{B^q} \\
  &= \frac 1 p \frac{\sum_k a_k^p}{A^p} + \frac 1 q \frac{\sum_k b_k^q}{B^q}
  = \frac 1 p + \frac 1 q = 1.
\end{align*}
And the theorem is proven. 

If $A=0$, all terms $a_k$ would be zero, and the inequality 
reduces trivially to $0\le 0$. The same holds if $B=0$.  
\end{proof}

In the case $p=q=2$, we obtain the 
\emph{Cauchy-Schwarz inequality}%
\mymargin{Cauchy-Schwarz inequality}%
\index{inequality!Cauchy-Schwarz}%
\index{inequality!Cauchy-Schwarz}%
\index{Cauchy-Schwarz!inequality}%
\index{Schwarz!Cauchy-Schwarz inequality}%
\[
   \sum_k \abs{a_k b_k} \le \sqrt{\sum_k a_k^2} \cdot \sqrt{\sum_k b_k^2}
\]
which, however, could be more easily proven 
for a generic inner product, as we will do 
in theorem~\ref{th:spazio_euclideo}.